\begin{definition}[Normalized characteristic-two refinement]
\label{mod:finitefield-chartworefinement}
Let $k$ be a finite field of characteristic two and put
\[
 \psi_0(x)=(-1)^{\operatorname{Tr}_{k/\mathbf F_2}(x)}.
\]
A normalized characteristic-two refinement is a Hasse function
$\mathfrak q:k\to\mathbf C^\times$ for the positive polar character
$\psi_0(xy)$ which is fixed by Frobenius.  Thus
\[
 \mathfrak q(x+y)=\mathfrak q(x)\mathfrak q(y)\psi_0(xy),
 \qquad
 \mathfrak q(x^2)=\mathfrak q(x).
\]
This is the Lean structure \texttt{CharTwoRefinement}.  Its values are
nonzero and unimodular, and Lean proves the exact identities
\[
 \mathfrak q(x)^2=\psi_0(x),
 \qquad
 \overline{\mathfrak q(x)}
   =\mathfrak q(x)\psi_0(x).
\]

Every Hasse function $h$ with this polar character can be normalized
by an affine twist.  More precisely, its Frobenius defect
$h(x^2)/h(x)$ is $\psi_0(cx)$ for a unique $c$.  One has
$\operatorname{Tr}(c)=0$, so
\[
 c=d+\operatorname{Frob}^{-1}(d)
\]
for some $d$, and $h(x)\psi_0(dx)$ is normalized.  Relative to any
fixed normalized $\mathfrak q$, every same-polar Hasse function is
uniquely
\[
 h_\lambda(x)=\mathfrak q(x)\psi_0(\lambda x).
\]
The affine member $h_\lambda$ is itself Frobenius-normalized exactly
when $\lambda=0$ or $\lambda=1$.  Hence the normalized refinements are
exactly $\mathfrak q$ and $\mathfrak q\psi_0$, and these two choices
are distinct.  Precomposition by Frobenius sends the affine coefficient
$\lambda$ to $\operatorname{Frob}^{-1}(\lambda)$.

The complete finite residual correction value on an odd critical line is
\[
 C_{\mathfrak q}(\lambda,c)
 =h_\lambda(c)
 =\mathfrak q(c)\psi_0(\lambda c)
 =\mathfrak q(\lambda)^{-1}\mathfrak q(c+\lambda).
\]
In particular it is $h_\lambda(c)$, not its inverse or conjugate.  If
\[
 S(\lambda)=\sum_{x\in k}h_\lambda(x),
 \qquad
 H(\lambda)=\Aphase(S(\lambda)),
 \qquad H_0=H(0),
\]
then $S(\lambda)\ne0$ and
$|S(\lambda)|=\sqrt{|k|}$.  Translation gives the correction rule
\[
 H(\lambda+c)=h_\lambda(c)^{-1}H(\lambda),
 \qquad
 h_\lambda(c)H(\lambda+c)=H(\lambda).
\]
Thus a later local linear phase must be transported together with this
complete positively oriented correction value, as required by
Lemma~\ref{lem:quadratic-correction-rule}; no local-field hypothesis is
assumed in this node.

The normalized phases satisfy
\[
 H(\lambda)=H_0\overline{\mathfrak q(\lambda)},
 \qquad
 H_0^2=\mathfrak q(1),
 \qquad
 H(\lambda)^2=\mathfrak q(1)\psi_0(\lambda),
\]
together with
\[
 H(\lambda^2)=H(\lambda),
 \qquad
 H(1+\lambda)H(\lambda)=1.
\]
Finally, Lean records the exact signed four-phase identity.  If
$c+d=a+b+A$, then
\[
 \frac{H(a)H(b)}{H(c)H(d)}
 =\mathfrak q(A)
  \psi_0\bigl((a+b)A+cd+ab\bigr).
\]
The numerator signs are $a,b$ and the denominator signs are $c,d$.
If also $B+C=A$, multiplication by the two positively oriented
refinement factors $\mathfrak q(B)\mathfrak q(C)$ leaves the phase
\[
 \psi_0\bigl((a+b)A+cd+ab+A+BC\bigr).
\]
Consequently, whenever
\[
 (a+b)A+cd+ab+A+BC+t=z+z^2,
\]
the Artin--Schreier identity gives the denominator-free cancellation
\[
 H(a)H(b)\mathfrak q(B)\mathfrak q(C)\psi_0(t)=H(c)H(d).
\]
This is purely finite-field algebra and assumes neither local break parity
nor the wild-quadratic conclusion.
\LeanFile{LanglandsFirstMainLemma/FiniteField/CharTwoRefinement.lean}
\lean{LanglandsFirstMainLemma.CharTwoRefinement}
\lean{LanglandsFirstMainLemma.CharTwoRefinement.map_add}
\lean{LanglandsFirstMainLemma.CharTwoRefinement.map_sq}
\lean{LanglandsFirstMainLemma.CharTwoRefinement.exists_charTwoRefinement_translate}
\lean{LanglandsFirstMainLemma.CharTwoRefinement.existsUnique_affine}
\lean{LanglandsFirstMainLemma.CharTwoRefinement.affine_frobenius}
\lean{LanglandsFirstMainLemma.CharTwoRefinement.affine_normalized_iff}
\lean{LanglandsFirstMainLemma.CharTwoRefinement.eq_or_eq_translateOne}
\lean{LanglandsFirstMainLemma.CharTwoRefinement.translateOne_ne}
\lean{LanglandsFirstMainLemma.CharTwoRefinement.correction_rule}
\lean{LanglandsFirstMainLemma.CharTwoRefinement.affineSum_ne_zero}
\lean{LanglandsFirstMainLemma.CharTwoRefinement.affineSum_norm}
\lean{LanglandsFirstMainLemma.CharTwoRefinement.correction_phase_rule}
\lean{LanglandsFirstMainLemma.CharTwoRefinement.value_sq}
\lean{LanglandsFirstMainLemma.CharTwoRefinement.value_conj}
\lean{LanglandsFirstMainLemma.CharTwoRefinement.affinePhase_eq}
\lean{LanglandsFirstMainLemma.CharTwoRefinement.affinePhase_zero_sq}
\lean{LanglandsFirstMainLemma.CharTwoRefinement.affinePhase_sq}
\lean{LanglandsFirstMainLemma.CharTwoRefinement.affinePhase_sq_arg}
\lean{LanglandsFirstMainLemma.CharTwoRefinement.affinePhase_pair}
\lean{LanglandsFirstMainLemma.CharTwoRefinement.affinePhase_four}
\lean{LanglandsFirstMainLemma.CharTwoRefinement.affinePhase_four_correction}
\lean{LanglandsFirstMainLemma.CharTwoRefinement.affinePhase_four_cancel}
\leanok
\SourceText{Lemmas~\ref{lem:normalized-refinement} and
\ref{lem:quadratic-correction-rule}.}
\uses{mod:finitefield-hassefunction,mod:finitefield-artinschreier}
\FormalizationContract
\end{definition}
\begin{proof}
For two same-polar Hasse functions, division cancels the polar factor
and produces an additive character.  Applying this to $h(x^2)$ and
$h(x)$ gives the Frobenius defect.  Its value at $1$ is one, hence its
coefficient has absolute trace zero.  Artin--Schreier exactness writes
that coefficient as $d+\operatorname{Frob}^{-1}(d)$; the adjoint
Frobenius identity then proves that twisting by $d$ is normalized.
Perfectness of the absolute-trace pairing proves uniqueness of every
affine coefficient.  Applying that uniqueness after Frobenius shows
that a normalized affine coefficient satisfies $\lambda^2=\lambda$,
so it is $0$ or $1$.

The affine translation and correction formulas are the Hasse functional
equation followed by finite-sum reindexing.  Nontriviality of $\psi_0$
and Hasse-sum orthogonality prove the asserted magnitude and nonvanishing
before any phase is simplified.  Setting the two arguments of the
functional equation equal gives $\mathfrak q(x)^2=\psi_0(x)$.
Conjugating the base sum identifies $H(1)$ with $\overline{H_0}$;
combining this with translation by $1$ gives
$H_0^2=\mathfrak q(1)$ and the remaining phase identities.

For the four-phase formula, first cancel the common factor $H_0^2$ in
the numerator and denominator.  The refinement law combines the four
remaining values.  Writing $Y=a+b$ and using $c+d=Y+A$, the apparent
extra exponent is $Y^2+Y$, whose normalized phase is one by the
Artin--Schreier identity.  A second use of the refinement law with
$B+C=A$ inserts the two refinement factors.  The final displayed
Artin--Schreier hypothesis then makes the residual phase exactly one.
\end{proof}
